% !TeX program = xelatex
\documentclass[UTF8,zihao=-4,fontset=fandol]{ctexart}
\newcommand{\ExamNumber}{3}
\usepackage{linear-algebra-exam}

\begin{document}
\ExamCover{7}{62}{50}

\Question{1}{单项选择}{2017 年数学一第 5 题（原卷第 1 页）}{4}{***}
设 $a$ 为 $n$ 维单位列向量，$E$ 为 $n$ 阶单位矩阵，则（\quad）。
\Choices
  {$E-aa^{\mathrm T}$ 不可逆。}
  {$E+aa^{\mathrm T}$ 不可逆。}
  {$E+2aa^{\mathrm T}$ 不可逆。}
  {$E-2aa^{\mathrm T}$ 不可逆。}

\Solution
\Known $a$ 是单位列向量，即各分量平方和为 $1$，因此 $a^{\mathrm T}a=1$。四个选项都在问某个矩阵是否不可逆。
\Method 四个矩阵都写成 $E+tP$，其中 $P=aa^{\mathrm T}$。由于中间两因子可合成 $a^{\mathrm T}a=1$，$P^2=P$。这意味着可以试用同样结构的 $E+sP$ 作为逆矩阵。
\Derive 把二者相乘，得
\[
(E+tP)(E+sP)=E+(t+s+ts)P.
\]
当 $t\ne-1$ 时，取 $s=-\frac{t}{1+t}$ 就能令括号中的系数为零。这个乘积换过顺序也为 $E$，因此 $E+P,E+2P,E-2P$ 都可逆。

当 $t=-1$ 时，$(E-P)a=a-a(a^{\mathrm T}a)=0$。若 $E-P$ 有逆矩阵，两边左乘该逆矩阵就会得出 $a=0$，这与单位列向量的条件矛盾。
先逐项检查 $P^2=(aa^{\mathrm T})(aa^{\mathrm T})=a(a^{\mathrm T}a)a^{\mathrm T}=P$，这里中间的 $1$ 是一个数。再把原选项的 $t$ 值代入 $s=-\frac{t}{1+t}$：当 $t=1,2,-2$ 时，$s$ 分别为 $-\frac{1}{2},-\frac{2}{3},-2$。这给出了 B、C、D 三个矩阵的具体逆矩阵。若 $t=-1$，分母为零，此时直接使用上面找到的非零列 $a$ 证明不可逆。
\Verify 选 \textbf{A}。例如 $E-2P$ 的平方为 $E-4P+4P^2=E$，因此它的逆矩阵正是它自身。

\Question{2}{单项选择}{2022 年数学一第 6 题（原卷第 1–2 页）}{4}{*****}
设 $A,B$ 为 $n$ 阶矩阵，$E$ 为单位矩阵。若方程组 $Ax=0$ 与 $Bx=0$ 同解，则（\quad）。

(A)\enspace 方程组 $\begin{pmatrix}A&O\\E&B\end{pmatrix}y=0$ 只有零解。

(B)\enspace 方程组 $\begin{pmatrix}E&A\\O&AB\end{pmatrix}y=0$ 只有零解。

(C)\enspace 方程组 $\begin{pmatrix}A&B\\O&B\end{pmatrix}y=0$ 与 $\begin{pmatrix}B&A\\O&A\end{pmatrix}y=0$ 同解。

(D)\enspace 方程组 $\begin{pmatrix}AB&B\\O&A\end{pmatrix}y=0$ 与 $\begin{pmatrix}BA&A\\O&B\end{pmatrix}y=0$ 同解。

\Solution
\Known $Ax=0$ 与 $Bx=0$ 同解，意味着对于任意 $n$ 维列向量 $x$，$Ax$ 为零当且仅当 $Bx$ 为零。各选项给出由 $A,B$ 组成的分块方程。
\Method 把 $y$ 分成两个 $n$ 维列向量 $u,v$，按分块乘法写出两个方程，比较对 $u,v$ 的限制。
\Derive C 的左式给出 $Au+Bv=0,\ Bv=0$，即 $Au=Bv=0$。右式同理等价于 $Bu=Av=0$。按题设，这两组条件相同。

A、B 可都取 $A=B=O$。无论哪个方程，此时都只限制 $u=0$，$v$ 任取。D 可取 $A=\begin{pmatrix}1&0\\0&0\end{pmatrix}$、$B=\begin{pmatrix}0&0\\1&0\end{pmatrix}$。二者乘同一列向量所得为零的条件都是其第一分量等于零，但 $AB=O$，$BA=B$。于是 $u=(1,0)^{\mathrm T},v=0$ 满足 D 的第一个方程，却使第二个方程的首块为 $Bu=(0,1)^{\mathrm T}$。
以 C 的左方程为例，从下块先读得 $Bv=0$，再代入上块 $Au+Bv=0$，得到 $Au=0$。反向代入这两个条件能恢复原方程，因此是等价改写。对右方程先读 $Av=0$，再得到 $Bu=0$，也可反向恢复。题设分别对 $u$ 和 $v$ 生效，这就足以比较两个方程的全部解。

D 的反例在两个乘积的顺序上产生差异。取题解中的 $u,v$ 时，D 两个方程的下块都为零；左式的上块是 $ABu+Bv=0$，右式的上块则是 $BAu+Av=(0,1)^{\mathrm T}$。因此确实有一个解只属于左方程。
\Verify 选 \textbf{C}。C 的两边都要求 $u,v$ 各自满足原题给出的同一组限制。

\Question{3}{单项选择}{2022 年数学二第 9 题（原卷第 2 页）}{4}{***}
设矩阵 $A=\begin{pmatrix}1&1&1\\1&a&a^2\\1&b&b^2\end{pmatrix}$，$\vec b=\begin{pmatrix}0\\2\\4\end{pmatrix}$，则线性方程组 $Ax=\vec b$ 的解的情况为（\quad）。
\Choices{无解。}{有解。}{有无穷多解或无解。}{有唯一解或无解。}

\Solution
\Known 系数矩阵的三行由 $1,a,b$ 分别生成，右侧三个数互不相同。要求判断所有参数情况下可能的解的种类。
\Method 先算行列式。系数行列式非零时可直接使用克拉默法则；行列式为零时，再看三行中哪两条方程发生冲突。
\Derive 从第二、三行各减第一行，按第一列展开，得
\[
|A|=\begin{vmatrix}a-1&a^2-1\\ b-1&b^2-1\end{vmatrix}
=(a-1)(b-1)(b-a).
\]
当 $1,a,b$ 两两不同时，$|A|\ne0$，故有唯一解。若其中两数相等，对应的两行系数完全相同，右侧数却不同。同一个左端不能同时取两个不同的值，故无解。
二阶展开可写成
\[
(a-1)(b^2-1)-(b-1)(a^2-1)
=(a-1)(b-1)\bigl[(b+1)-(a+1)\bigr].
\]
因此只需检查 $a=1$、$b=1$、$a=b$ 三种退化条件。第一种使第一、二式左端相同，右端是 $0,2$；第二种使第一、三式左端相同，右端是 $0,4$；第三种使第二、三式左端相同，右端是 $2,4$。三种情况都产生矛盾，重叠时也一样。
\Verify 选 \textbf{D}。例如 $(a,b)=(2,3)$ 给出唯一解，$a=b=2$ 时第二、三条方程的左侧相同而右侧分别是 $2$ 和 $4$。

\Question{4}{填空}{2015 年数学一第 13 题（原卷第 2 页）}{10}{****}
$n$ 阶行列式
\[
\begin{vmatrix}
 2 & 0 & \cdots & 0 & 2\\
-1 & 2 & \cdots & 0 & 2\\
\vdots & \vdots & \ddots & \vdots & \vdots\\
0 & 0 & \cdots & 2 & 2\\
0 & 0 & \cdots & -1 & 2
\end{vmatrix}
=\blank
\]

\Solution
\Known 记题中的行列式为 $D_n$。主对角线和末列都是 $2$，主对角线正下方是 $-1$，其余位置为零。要求用 $n$ 表示 $D_n$。
\Method 当 $n\ge3$ 时，首列只有前两项非零。按首列展开，第一项的余子式仍有相同结构，从而可把问题化为低一阶的行列式。
\Derive 首列第一项为 $2$，其贡献是 $2D_{n-1}$。第二项是 $-1$，代数余子式的符号是 $(-1)^{2+1}=-1$，两个负号相乘为正。对应余子式划掉原第 $2$ 行、第 $1$ 列，剩下的第一行只有末项 $2$。再按这一行展开，得到 $2(-1)^n$ 乘以一个 $n-2$ 阶上三角行列式，其对角线都是原次对角线上的 $-1$。因此第二项的贡献是
\[
2(-1)^n(-1)^{n-2}=2,
\qquad D_n=2D_{n-1}+2.
\]
直接算得 $D_2=\begin{vmatrix}2&2\\-1&2\end{vmatrix}=6$。将递推写成 $D_n+2=2(D_{n-1}+2)$，得 $D_n+2=2^{n-2}(6+2)=2^{n+1}$。
余子式中第二次展开时，末项所在位置为第 $1$ 行、第 $n-1$ 列，符号是 $(-1)^{1+(n-1)}$。删去这一行、这一列后，原第 $3,\ldots,n$ 行和第 $2,\ldots,n-1$ 列对齐，每行最靠左的非零元正是 $-1$，故剩下的上三角行列式值为 $(-1)^{n-2}$。例如 $n=3$ 时递推给出 $D_3=2\cdot6+2=14$，也可直接沿第一列展开获得。
\Verify 填 $\boxed{2^{n+1}-2}$。代入 $n=2$ 得到 $6$，和直接计算相同。若把该结构延伸至 $n=1$，公式也给出 $D_1=2$。

\Question{5}{填空}{2021 年数学一第 15 题（原卷第 2 页）}{10}{**}
设 $A=(a_{ij})$ 为 $3$ 阶矩阵，$A_{ij}$ 为元素 $a_{ij}$ 的代数余子式，若 $A$ 的每行元素之和均为 $2$，且 $|A|=3$，则 $A_{11}+A_{21}+A_{31}=\blank$。

\Solution
\Known 左边是同一列中三个元素的代数余子式之和。已知 $A$ 每行元素的和为 $2$，且行列式为 $3$。
\Method 按第一列展开时，这三个代数余子式恰好与第一列的元素分别相乘。若能把第一列全变成 $2$，所求的和就会出现。代数余子式在计算时要删去第一列，所以替换第一列的元素不会改变它们。
\Derive 将第二列、第三列分别加到第一列，行列式保持为 $3$。新的第一列的每个元素都是对应行的元素和，也就是 $2$。按新第一列展开得
\[
3=2A_{11}+2A_{21}+2A_{31}\text{。}
\]
把替换第一列分成两次列加法即可说明行列式不变。先加第二列，再加第三列，每次都把另一列的一倍加到第一列。三个目标余子式删去的都是第一列，剩下两列完全没有变化，位置符号也仍是 $+,-,+$。因此展开新第一列时可以继续写原题的 $A_{11},A_{21},A_{31}$。
\Verify 填 $\boxed{\frac{3}{2}}$。可用伴随矩阵再检查一次。令 $u=(1,1,1)^{\mathrm T}$，则 $Au=2u$。左乘 $A^*$ 得 $3u=2A^*u$，取第一分量即得同一结果。

\Question{6}{填空}{2023 年数学二第 16 题（原卷第 2 页）}{10}{****}
已知线性方程组
\[
\begin{cases}
ax_1+x_3=1,\\
x_1+ax_2+x_3=0,\\
x_1+2x_2+ax_3=0,\\
ax_1+bx_2=2
\end{cases}
\]
有解，其中 $a,b$ 为常数。若
$\begin{vmatrix}a&0&1\\1&a&1\\1&2&a\end{vmatrix}=4$，
则 $\begin{vmatrix}1&a&1\\1&2&a\\a&b&0\end{vmatrix}=\blank$。

\Solution
\Known 四条方程有一组共同的解。已知前三条方程的系数行列式为 $4$，要求后三条方程的系数行列式。
\Method 把前三列系数和等号右侧的常数排成四列。题设保证有解，因此第四列是前三列按解的分量加权后相加的结果。用行列式的列加法把第四列变为零列，就能把两个三阶行列式联系起来。
\Derive 设共同的解为 $(u,v,w)$，四列分别记为 $C_1,C_2,C_3,C_4$。四条方程恰好说明 $C_4=uC_1+vC_2+wC_3$。第四列依次减去前三列的 $u,v,w$ 倍后变为零列，行列式的值保持不变。所以
\[
0=\begin{vmatrix}
a&0&1&1\\
1&a&1&0\\
1&2&a&0\\
a&b&0&2
\end{vmatrix}
=-\begin{vmatrix}1&a&1\\1&2&a\\a&b&0\end{vmatrix}
 +2\begin{vmatrix}a&0&1\\1&a&1\\1&2&a\end{vmatrix}.
\]
这里按第四列展开，位置 $(1,4)$ 的符号为 $(-1)^5=-1$，位置 $(4,4)$ 的符号为 $(-1)^8=1$。
按第四列展开时，第二、第三项因常数为零而消失。删第一行、第四列留下的是后三式的系数矩阵；删第四行、第四列留下的是前三式的系数矩阵。设前者行列式为 $D$，则等式恰为 $0=-D+2\cdot4$。例如取 $a=2,b=-4$，四式的共同解为 $(\frac{1}{2},-\frac{1}{4},0)$，第四式左端为 $1+1=2$，两组系数行列式分别为 $4,8$。
\Verify 填 $\boxed{8}$。代回第四列展开式为 $-8+2\cdot4=0$，符合四条方程有共同解时得到的零行列式。

\Question{7}{解答}{2015 年数学二第 22 题（原卷第 4 页）}{20}{****}
设矩阵 $A=\begin{pmatrix}a&1&0\\1&a&-1\\0&1&a\end{pmatrix}$，且 $A^3=O$。

（I）求 $a$ 的值。

（II）若矩阵 $X$ 满足 $X-XA^2-AX+AXA^2=E$，其中 $E$ 为 $3$ 阶单位矩阵，求 $X$。

\Solution
\Known 第一问从 $A^3=O$ 求出 $A$ 中的参数。第二问从含四项的矩阵等式中解出 $X$，其中 $O$ 是零矩阵。
\Method 第一问对 $A^3$ 取行列式，得到关于 $a$ 的方程。第二问按 $X$ 左、右两侧的因子分别整理。由于三次及更高次的 $A$ 的幂都为零，逆矩阵可以通过有限次乘法直接得到。
\Derive \textbf{（I）} $0=|A^3|=|A|^3$，所以 $|A|=0$。按第一行展开，
\[
|A|=a(a^2+1)-a=a^3\text{，}
\]
故 $a=0$。把此值代入原矩阵，可算出 $A^2=\begin{pmatrix}1&0&-1\\0&0&0\\1&0&-1\end{pmatrix}$，再右乘 $A$ 得 $O$。所以 $a=0$ 确实满足题设。

\textbf{（II）} 原式先在右侧提取 $E-A^2$，成为 $X(E-A^2)-AX(E-A^2)$，再在左侧提取 $E-A$，得到 $(E-A)X(E-A^2)=E$。用乘法展开检查两个系数的逆：
\[
(E-A)(E+A+A^2)=E-A^3,\qquad
(E-A^2)(E+A^2)=E-A^4.
\]
题设给出 $A^3=O$，所以 $A^4=AA^3=O$。两式右侧都等于 $E$。两式中的各因子都是 $A$ 的多项式，颠倒乘法顺序仍得到 $E$。因此 $E-A$ 和 $E-A^2$ 都可逆。
\SolutionContinued

在 $(E-A)X(E-A^2)=E$ 两侧同左乘 $(E-A)^{-1}$、同右乘 $(E-A^2)^{-1}$，恰好消去 $X$ 两边的因子，得到唯一解
\[
X=(E+A+A^2)(E+A^2)=E+A+2A^2\text{。}
\]
展开解的乘积时，留下 $E$、一个 $A$ 和两个 $A^2$，含三次、四次幂的项全为零。最后代入第一问的 $A$ 和 $A^2$，各位置分别相加即可。
\Verify $a=0$，$X=\begin{pmatrix}3&1&-2\\1&1&-1\\2&1&-1\end{pmatrix}$。该 $A$ 满足 $A^2A=O$。由 $X=E+A+2A^2$ 得 $(E-A)X=E+A^2$，再右乘 $E-A^2$ 确为 $E$。
\end{document}
